%! TEX root = part1.tex
% the main TeX file which is intended to compile, :VimtexReload after adjustment
   % this file should be included in the main file
% :h vimtex-tex-root
\documentclass[../main.tex]{subfiles}
% ensure the class of docs
% enumerate may not work

\begin{document}
\begin{frame}[t,mycolor=digiPH_white,mytitle=standard,light]
	\frametitle{Pengertian Arsitektur \& Struktur}

	\textbf{Arsitektur}
	\tabitem Seni bangunan yang telah ada sejak manusia ada
	\tabitem Disebut seni terikat karena:
	\begin{itemize}
		\item Digunakan manusia
		\item Material memiliki keterbatasan teknis
	\end{itemize}
	\tabitem \textit{Memadukan:}
	\begin{itemize}
		\item Falsafah
		\item Religi
		\item Tradisi
		\item Seni
		\item Ilmu pengetahuan
	\end{itemize}

\end{frame}


\begin{frame}[t,mycolor=digiPH_red,mytitle=standard,dark]

	\frametitle{Fungsi Struktur Bangunan}

	\begin{itemize}
		\item Menyalurkan beban ke tanah
		\item Menahan:
		      \begin{itemize}
			      \item Beban mati
			      \item Beban hidup
			      \item Beban angin
			      \item Beban gempa
		      \end{itemize}
		\item Menjamin keamanan dan stabilitas bangunan
	\end{itemize}

	\footnotesize Struktur = sistem yang membawa beban dari atas → ke bawah → ke tanah melalui pondasi.
\end{frame}


\begin{frame}[t,mycolor=digiPH_red,mytitle=standard,dark]
	\frametitle{Definisi Sistem Struktur}

	\txt{digiPH_ocean}{digiPH_white}{Struktur Bangunan Gedung adalah organisasi elemen-elemen bangunan yang mendukung berfungsinya bangunan dengan baik}


\end{frame}

\begin{frame}[t,mycolor=digiPH_red,mytitle=standard,dark]
	\frametitle{Sistem Struktur}

	\tabitem Bentuk organisasi elemen struktur
	\tabitem Bertujuan menyalurkan beban secara karakteristik


\end{frame}
\begin{frame}[t,mycolor=digiPH_leaf,mytitle=standard,dark]
	\frametitle{Contoh Elemen Struktur:}
	\begin{itemize}
		\item Pondasi
		\item Sloof
		\item Kolom
		\item Balok
		\item Ring balok
		\item Kuda-kuda
		\item Atap
	\end{itemize}

\end{frame}

\begin{frame}[t,mycolor=digiPH_leaf,mytitle=standard,dark]


	\frametitle{Prinsip Kerja Struktur}

	\begin{itemize}
		\item Meneruskan beban bangunan dari atas ke bawah
		\item Menyebarkan beban ke tanah
		\item Menjamin:
		      \begin{itemize}
			      \item Kekuatan
			      \item Stabilitas
			      \item Keamanan
		      \end{itemize}
		\item Struktur harus mampu menahan:
		      \begin{itemize}
			      \item Beban vertikal (mati \& hidup)
			      \item Beban horizontal (angin \& gempa)
		      \end{itemize}
	\end{itemize}

\end{frame}

\begin{frame}[c,mycolor=digiPH_white,mytitle=standard,light]
	\frametitle{Jenis Penyaluran Beban}
	\tabitem Struktur Utama
	\begin{itemize}
		\item Menyalurkan beban ke tanah
		\item Tanpa struktur utama → bangunan tidak dapat berfungsi
	\end{itemize}
	\tabitem Struktur Pendukung
	\begin{itemize}
		\item Mendukung kerja struktur utama
		\item Membantu distribusi beban
	\end{itemize}

\end{frame}


\begin{frame}[c,mycolor=digiPH_gray,mytitle=standard,light]

	\frametitle{Jenis-Jenis Beban Struktur}

	\ltxt{digiPH_gray}{digiPH_black}{\tabitem  \textbf{Beban Dasar}
		\begin{itemize}
			\item Beban Mati
			\item Beban Hidup
			\item Beban Angin
			\item Beban Termis
			\item Gerakan tanah
			\item Gempa bumi
			\item Beban Dinamis
		\end{itemize}}

	\rtxt{digiPH_gray}{digiPH_black}{\tabitem  \textbf{Beban Tambahan}
		\begin{itemize}
			\item Beban Salju
			\item Beban Hujan
			\item Beban Ledakan
			\item Beban Tanah
			\item Beban Air
		\end{itemize}}


\end{frame}

\begin{frame}[t,mycolor=digiPH_leaf,mytitle=standard,dark]

	\frametitle{Peran Struktur Bangunan}
	\utxt{digiPH_gray}{digiPH_black}{Estetika \\
		Memberikan keindahan dan keserasian}

	\btxt{digiPH_gray}{digiPH_black}{Fungsional \\
		Memberikan kenyamanan dan kemanfaatan}

\end{frame}



\begin{frame}[t,mycolor=digiPH_leaf,mytitle=standard,dark]
	\frametitle{Peran Struktur Bangunan}
	\utxt{digiPH_gray}{digiPH_black}{Struktural \\
		Kuat dan stabil}

	\btxt{digiPH_gray}{digiPH_black}{Ekonomis \\
		Efisien, proporsional, tahan lama}
\end{frame}

\begin{frame}[c,mycolor=digiPH_white,mytitle=standard,light]

	\begin{table}[h]
		\centering
		\caption{Perkembangan Arsitektur Berdasarkan Peradaban}
		\begin{tabular}{|l|p{3cm}|p{3cm}|}
			\hline
			Peradaban          &  Contoh Bangunan                 &  Karakteristik  \\
			\hline
			Prasejarah         &  Gua, gubuk                      &  Perlindungan dasar  \\
			Kuno               &  Piramida, Parthenon, Colosseum  &  Monumental \& religius  \\
			Abad Pertengahan   &  Katedral Gotik                  &  Lengkung runcing, kaca patri  \\
			Revolusi Industri  &  Jembatan, pabrik                &  Material besi \& baja  \\
			Modernisme         &  Pencakar langit                 &  Baja \& kaca  \\
			Kontemporer        &  Gedung pintar                   &  Berkelanjutan \& efisiensi energi  \\
			\hline
		\end{tabular}
	\end{table}
\end{frame}


\begin{frame}[t,mycolor=digiPH_ocean,mytitle=center,dark]
	\frametitle{Dasar Konstruksi Bangunan}
	\begin{itemize}
		\item Prinsip kekuatan dan stabilitas
		\item Efisiensi material
		\item Keberlanjutan
		\item Keamanan
		\item Manajemen proyek
	\end{itemize}
\end{frame}

\begin{frame}[t,mycolor=digiPH_ocean,mytitle=center,dark]

	\textbf{Struktur harus mampu menahan:}

	\begin{itemize}
		\item Beban vertikal
		\item Beban horizontal
	\end{itemize}
\end{frame}

\begin{frame}[t,mycolor=digiPH_white,mytitle=standard,light]

	\tabitem Sub Structure \textit{(Struktur Bawah)}
	\begin{itemize}
		\item Terletak di bawah tanah
		\item Contoh: Pondasi
		\item Fungsi: Menerima dan menyalurkan gaya ke tanah
	\end{itemize}

	\tabitem Middle Structure \textit{(Struktur Tengah)}
	\begin{itemize}
		\item Di atas tanah, di bawah atap
		\item Kolom, balok, pelat lantai
		\item Menyalurkan gaya dalam bangunan
	\end{itemize}

	\tabitem Super/Up Structure \textit{(Struktur Atas)}
	\begin{itemize}
		\item Menopang atap
		\item Contoh: Kuda-kuda
		\item Menyalurkan beban atap
	\end{itemize}
\end{frame}

\begin{frame}[t,mycolor=digiPH_leaf,mytitle=standard,dark]

	\tabitem Sistem konstruksi:
	\begin{itemize}
		\item Elemen yang menempel pada struktur utama
		\item Berfungsi menyebarkan gaya dan menerima beban langsung
	\end{itemize}

	\textbf{Penempatannya:}
	\tabitem Super Struktur
	\begin{itemize}
		\item Tangga
		\item Dinding
		\item Plafond
	\end{itemize}

	\tabitem Up Struktur
	\begin{itemize}
		\item Atap
		\item Listplank
		\item Talang air
	\end{itemize}
\end{frame}

\begin{frame}[t,mycolor=digiPH_ocean,mytitle=center,dark]

	Struktur adalah sistem penyalur beban dari bangunan ke tanah.

	\only<1>{\tabitem Struktur menentukan:
		\begin{itemize}
			\item Keamanan
			\item Kenyamanan
			\item Umur bangunan
		\end{itemize}}

	\only<2>{\tabitem Perancangan struktur harus mempertimbangkan:
		\begin{itemize}
			\item Beban
			\item Stabilitas
			\item Efisiensi
			\item Keberlanjutan
		\end{itemize}}

\end{frame}

\begin{frame}[c,mycolor=digiPH_leaf,mytitle=standard,dark]
	\frametitle{Tugas Kelas:}

	Tuliskan ringkasan (resume) dari literatur terkait elemen-elemen struktur. maks. 2 halaman. ;

\end{frame}

\begin{comment}

\begin{frame}[c,mycolor=digiPH_gray,mytitle=standard,light]
	\frametitle{Tugas Rumah}

	% \bgimg{1}{denahrmh}

\end{frame}

\begin{frame}[t,mycolor=digiPH_gray,mytitle=standard,light]
	\frametitle{Minggu Depan}

	Tolong Bawa:
	\begin{itemize}
		\item Kertas Karton
		\item Tusuk sate
		\item Alat pembuatan maket (gunting, cutter, lem fox, lem uhu, lem korea, dll)
	\end{itemize}

\end{frame}

\end{comment}




\bibliographystyle{apalike}
\bibliography{../biblio.bib} % required for citecompletion, not work in mainfile
\end{document}

